\documentclass{article}

% arXiv preprint in NeurIPS style. The [preprint] option produces a
% non-anonymous, single-column preprint suitable for arXiv upload.
% neurips_2024.sty is bundled alongside this file.
\usepackage[preprint]{neurips_2024}

\usepackage[utf8]{inputenc}
\usepackage[T1]{fontenc}
\usepackage{hyperref}
\usepackage{url}
\usepackage{booktabs}
\usepackage{amsmath,amssymb,amsthm}
\usepackage{amsfonts}
\usepackage{microtype}
\usepackage{xcolor}

\newcommand{\R}{\mathbb{R}}
\newcommand{\Hyp}{\mathbb{H}}
\newcommand{\kp}{\kappa}
\newcommand{\code}[1]{\texttt{#1}}

\title{Meaning on a Video Card:\\ Scalable, Lossless Geometric Primitives}

% ORCID and contact are placeholders; fill before upload.
\author{%
  John Vaught \\
  Synoros \\
  \texttt{ORCID: <fill-in before upload>}%
}

\begin{document}

\maketitle

\begin{abstract}
\code{holonomy\_lib} is a PyTorch library of geometric and spectral
primitives --- Riemannian manifolds, hyperbolic graph operations,
spectral graph theory, discrete Ricci curvature, computational topology
--- built as a trusted substrate for learning geometric structure
directly. Seven of its primitives surfaced results worth recording, each
from the work of making the mathematics hold under \code{float64},
automatic differentiation, and GPU execution: (i) an
automatic-differentiation pitfall in the hyperbolic
heat-kernel dimension recursion, where a multiplicative spectral-shift
factor is silently dropped; (ii) a hand-derived closed form for the
$\Hyp^5$ heat kernel that is faster and roughly three orders of
magnitude more precise than the recursion; (iii) an arcsinh
reparameterization of hyperbolic distance that removes two distinct
backward-\code{NaN} modes without an $\varepsilon$ hyperparameter;
(iv) a curvature-sign dispatch that lets a learnable curvature $\kp$
cross zero during training without re-instantiation; (v) a continuous
per-point-$\kp$ manifold that improves on discrete curvature gating in a
controlled identifiability task; (vi) a pair-$\kp$ combiner abstraction
with an empirical characterization of the design axis; and (vii) a
reference-free dual-check validation method for operator-semigroup
primitives, which caught (i). Each result lives in the gap between a
formula correct on paper and code that holds to it in \code{float64}
under autograd, at practical speed on commodity GPUs. Two further
capabilities study the geometry once it is correct: a spectral-dimension
estimator that detects representational collapse, and a
content-addressable provenance layer that makes activation tracing and
ablation native operations. Together they make one case --- that
nonlinear, non-Euclidean embedding is practically viable, not only
theoretically sound. Section~\ref{sec:intro} places the library in the
research program it serves.
\end{abstract}

\section{Introduction}
\label{sec:intro}

\begin{quote}
\itshape
When I first heard Sutton say that large language models are bitter too,
it crystallized a realization I had been circling in theory: that what a
language model learns is the geometry of meaning --- the shape of how
things relate --- and that it is made to learn this only indirectly, as
an artifact of compressing how we communicate, rather than from the shape
itself. \upshape\textemdash\,J.~Vaught
\end{quote}

\citet{sutton2019bitter} argued that across six decades of artificial
intelligence, general methods that scale with computation --- search and
learning --- have consistently overtaken systems built from
human-encoded domain knowledge. The knowledge-engineering approach wins
in the short run and is overtaken in the long run; this is the bitter
lesson. In a 2025 interview \citep{sutton2025dwarkesh}, Sutton extended
the argument to current large language models: they learn by imitating a
finite store of human-generated text and do not learn continually from
their own experience, so he expects systems that learn from experience
and computation to eventually supersede them --- another turn of the
same lesson, with the models that depend most on human-produced data
being the ones overtaken.

One reading of this is structural. Language is itself a human artifact:
an engineered, lossy intermediate representation of the world. Models
trained on it do acquire non-linguistic structure --- linear analogy
geometry in word embeddings \citep{mikolov2013linguistic}, the
Fourier-basis representation a small transformer forms to do modular
addition \citep{nanda2023progress}, and the geometric addition circuit a
production model uses in place of the textbook algorithm
\citep{lindsey2025biology}. But that structure is acquired
\emph{indirectly}: reconstructed from the statistics of text rather than
operated on directly, and at a cost in data and compute far above what
the structure itself requires. If the bitter lesson applies one level
down, the language channel is the human-encoded prior that scale will
eventually route around.

The compute that powers that scale is no more neutral. In practice scale
means the GPU, and the GPU is a human-engineered artifact built to
rasterize triangles and decode video, not to represent meaning. The
general methods that overtake hand-built knowledge run on hardware
hand-built for graphics --- dense regular arithmetic at high throughput,
where exact, lossless representation of structure was never a design
goal. The substrate the bitter lesson runs on is itself a domain-specific
human prior. Geometric machinery that means to ride the same hardware
therefore carries a concrete obligation: stay numerically exact while
fitting what a video card does well.

Synoros was founded on the wager that this verdict can be acted on
rather than accepted: that geometric and algorithmic structure can be
made \emph{directly available} to learning --- as differentiable,
composable primitives a model optimizes over and through --- instead of
rediscovered through language. The working hypothesis of this substrate
research program is that supplying correct geometric machinery, and
letting scale and search act on it directly, reaches the structure that
language models acquire slowly and as a side effect by a shorter and
less wasteful route. The program is pursued in a sibling
project (\code{synoros-substrate}) and shares the direction of an active
literature on curvature-aware representation learning
\citep{nickel2017poincare,gu2019learning,bachmann2020constant,digiovanni2022heterogeneous,guo2025graphmore}.
A learner that operates on a geometric substrate, however, inherits
every error in that substrate. The substrate has to be correct first.

The case for curved geometry here is target-fidelity, not analogy
\citep{vaught2026afteruniverse}. The familiar version --- spacetime is
hyperbolic, so embeddings should be hyperbolic --- is loose enough that a
careful reader dismisses it, and rightly. The stronger claim is that modeling meaning is, with one label
peeled away, modeling the world that meaning describes, and an honest
model must span the regimes that world presents rather than the single
flat regime a linear-algebra default supplies. Aircraft engineers studied
birds without claiming aircraft are birds: they took the features that
make flight work and built systems that respect them. The analogue here
is not a curved manifold because the universe is curved, but a substrate
that spans curvature --- spherical, flat, hyperbolic, and the per-point
mixtures and sign changes between them --- because a faithful model of
the target cannot be confined to one. The target is the structure of
relations, not sensory data: it can be learned from text, images, or
concepts alike, since what is fit is the shape of those relations and not
the semantics that express them.

\code{holonomy\_lib} is that substrate layer. It consolidates the
mathematics that geometric machine learning otherwise re-implements
project by project: Riemannian manifolds (fixed-rank, SPD, the
hyperboloid model, the $\kp$-stereographic model, the Lorentzian
$(1,n{-}1)$ model, product manifolds, and a per-point-$\kp$
heterogeneous manifold); hyperbolic graph operations (the Fr\'echet mean
\citep{karcher1977riemannian}, manifold-aware inner products
\citep{pennec2006intrinsic}, hyperbolic Laplacian eigenmaps, the
$\Hyp^n$ heat kernel); spectral graph theory (combinatorial, normalized,
signed and magnetic Laplacians, diffusion maps, effective resistance,
spectral dimension); discrete Ricci curvature and flow; tensor
decompositions; Riemannian optimization \citep{absil2008optimization};
simplicial topology and batched
persistent homology; cellular sheaves; \mbox{SO(3)} primitives; and a
content-addressable provenance layer. Everything is GPU-native and
batched-first; every numerical constant is either derived from inputs,
marked as a universal invariant, or cataloged with a documented
procedure and scale of validity; every public primitive carries a
citation to the source that defines its mathematics; and a static audit
enforces the constant discipline in continuous integration.

Optimizing that substrate took work in seven specific places, each
between mathematical correctness and implementation correctness under
\code{float64}, automatic differentiation, and GPU execution. The claim
is narrow and concrete: nonlinear, non-Euclidean embedding --- learnable
curvature that is per-point and free to cross zero --- is practically
viable, not only theoretically sound, because the primitives stay
numerically exact and run at practical speed on commodity hardware (an
AMD RX 9060 XT under ROCm, an NVIDIA RTX 3060 under CUDA; no datacenter
accelerator). A substrate earns the trust to build on only once these
gaps close, and the sections that follow close them: correctness
(Section~\ref{sec:correctness}), precision and speed
(Section~\ref{sec:precision}), new primitives
(Section~\ref{sec:primitives}), and a provenance layer that keeps the
substrate auditable (Section~\ref{sec:provenance}).
Section~\ref{sec:discussion} returns to the framing.

\paragraph{Artifacts.} Every result below is backed by a symbolic
verification script, a numerical demonstration, and a consolidated
write-up in the repository under \code{notes/verification/},
\code{notes/strengthening/}, and \code{notes/validation/}. Specific
paths appear inline.

\section{Background}
\label{sec:background}

\paragraph{The $\kp$-stereographic model.} A single model of constant
sectional curvature $\kp \in \R$ interpolates the sphere ($\kp > 0$),
Euclidean space ($\kp = 0$), and hyperbolic space ($\kp < 0$). Distances
and the exponential/logarithm maps are expressed through the
curvature-dependent trigonometric functions
$\tan_\kp$ and its inverse $\tan_\kp^{-1}$, built on M\"obius gyrovector
operations \citep{ungar2008gyrovector}; \citet{bachmann2020constant}
introduced this model to graph convolutional networks with $\kp$ a
single learnable scalar per layer, recovering the Euclidean network as
$\kp \to 0$. Mixed-curvature representations take products of such spaces
with fixed per-factor curvatures \citep{gu2019learning,skopek2020mixed}.

\paragraph{The hyperbolic heat kernel.} The heat kernel on $\Hyp^n$ is a
radial function $k_t(r)$ solving the heat equation
$\partial_t k = \Delta k$ with a point-mass initial condition
\citep{davies1989heat}.
\citet{daviesmandouvalos1988} give the $\Hyp^3$ closed form; across
dimension, the kernel obeys a recursion that raises $n$ by two by
applying the operator $\frac{1}{\sinh r}\,\partial_r$
\citep{grigoryan1998heat,grigoryan2009heat,yu2018heat}; both appear in
Sections~\ref{sec:c1} and~\ref{sec:c2}.

\paragraph{Per-point curvature.} Beyond a single learnable $\kp$, several
constructions allow curvature to vary: products with spherically
symmetric factors \citep{digiovanni2022heterogeneous}, reinforcement-learning
control of a global $\kp$ \citep{fu2021ace}, and --- closest to the
per-node setting --- a mixture of Riemannian experts that gates each node
over a \emph{discrete} bank of constant-curvature spaces
\citep{guo2025graphmore}. Section~\ref{sec:c5} contrasts a continuous
per-point $\kp$ against this discrete gating.

\paragraph{Numerical setting.} All primitives are implemented on
\code{torch.Tensor} with automatic differentiation
\citep{paszke2019pytorch}. Where an established library exists, the
implementation is checked against it: \code{geoopt} for Riemannian operations
\citep{kochurov2020geoopt}, whose $\varepsilon$-clamp idiom appears in
Section~\ref{sec:c3}, and the \code{safe\_*} helper pattern shared with
PyTorch3D \citep{ravi2020pytorch3d}.

\section{Implementation correctness}
\label{sec:correctness}

\subsection{An autograd pitfall in the heat-kernel dimension recursion}
\label{sec:c1}

The recursion that raises the hyperbolic heat kernel by two dimensions,
stated correctly in the literature
\citep{grigoryan1998heat,grigoryan2009heat,yu2018heat}, is
\begin{equation}
  k^{\,n+2}(t, r) \;=\; -\,\frac{e^{-n t}}{2\pi \sinh r}\,
  \partial_r\, k^{\,n}(t, r).
  \label{eq:recursion}
\end{equation}
The factor $e^{-n t}$ is a spectral shift: it is multiplicative and
\emph{constant in $r$}. This is exactly what makes it easy to drop when
porting~\eqref{eq:recursion} to automatic differentiation. The natural
implementation,
\begin{center}
\code{grad = torch.autograd.grad(k\_n(t, r).sum(), r)[0];}\quad
\code{return -grad / (2*math.pi*torch.sinh(r))},
\end{center}
computes $\partial_r k^{\,n}$ faithfully and then divides by
$2\pi \sinh r$ --- but never multiplies by $e^{-nt}$, because that factor
does not depend on the differentiated variable $r$ and so is invisible
to \code{torch.autograd.grad}. The factor has to be applied by hand,
outside the autograd call. The corrected kernel multiplies by
$e^{-n_{\text{prev}} t}$ explicitly and clamps $\sinh r$ at
\code{finfo.tiny} for the $r \to 0$ limit.

The error has a clean signature. At a sample $(t, r) = (0.5, 1.0)$ the
naive recursion returns exactly $e^{3 \cdot 0.5} = e^{1.5} \approx 4.4817$
times the correct value --- the missing $e^{nt}$. Symbolic verification
(\code{notes/verification/heat\_kernel\_recursion\_sympy.py}) confirms
this algebraically: \code{sympy.simplify} reduces the heat-equation
residual of the corrected $k^5$ to exact zero, while the naive $k^5$
leaves a nonzero residual.

The generalizable lesson: when porting an analytical formula to autograd
code, any multiplicative factor that depends on a \emph{non-differentiated}
variable --- here $t$ and the dimension $n$ --- must be applied by hand.
Automatic differentiation reports only the derivative it computes, so a
factor outside that derivative stays silently absent.

\subsection{Equivalent mathematics, very different float behavior:
arcsinh vs.\ arccosh}
\label{sec:c3}

The textbook curvature-$\kp$ hyperbolic distance
$d_\kp(x,y) = \tfrac{1}{\sqrt{|\kp|}} \operatorname{arccosh}
\!\big(\kp \langle x, y\rangle_{\mathcal{M}}\big)$ is exact on paper and
unreliable in \code{float64} under autograd, for two distinct reasons:
\begin{enumerate}\itemsep2pt
  \item \emph{Boundary derivative singularity.}
  $\frac{d}{dz}\operatorname{arccosh}(z) = 1/\sqrt{z^2-1}$ diverges at
  $z = 1$, i.e.\ at the on-manifold self-pair $x = y$. The forward value
  $\operatorname{arccosh}(1) = 0$ is finite, but the backward pass
  propagates $\infty \cdot \partial z/\partial x = \code{NaN}$.
  \item \emph{Catastrophic cancellation.} For $x \approx y$,
  $z = \kp\langle x,y\rangle_{\mathcal M} = \cosh(\alpha) = 1 + \alpha^2/2 +
  O(\alpha^4)$, so the information about a small distance $\alpha$ sits in
  the last bits of a number near $1$; $\operatorname{arccosh}(1+\delta)
  \approx \sqrt{2\delta}$ then amplifies the lost precision.
\end{enumerate}
The equivalent reparameterization
\begin{equation}
  d_\kp(x, y) \;=\; \frac{2}{\sqrt{|\kp|}}\,
  \operatorname{arcsinh}\!\Big(\tfrac{1}{2}\sqrt{|\kp|}\;
  \lVert y - x \rVert_{\mathcal M}\Big)
  \label{eq:arcsinh}
\end{equation}
removes both. It computes $\lVert y - x\rVert_{\mathcal M}^2$ from
coordinate differences (no subtraction near $1$), and $\operatorname{arcsinh}$
is entire, with derivative $1/\sqrt{1+\arg^2}$ equal to $1$ at the origin
--- no boundary singularity. The half-angle identity
$\cosh A - 1 = 2\sinh^2(A/2)$ makes~\eqref{eq:arcsinh} algebraically
identical to the textbook form
(\code{notes/verification/arcsinh\_reparam\_sympy.py}).
Table~\ref{tab:arcsinh} shows the practical difference in \code{float64}.

\begin{table}[t]
  \centering
  \caption{Recovered distance at small geodesic separation $\alpha$,
  three implementations. The textbook form loses all bits below
  $\alpha \sim 10^{-8}$; the $\varepsilon$-clamp
  (\code{geoopt}-style, $\varepsilon = 10^{-5}$) returns the constant
  $\operatorname{arccosh}(1+\varepsilon) \approx \sqrt{2\varepsilon}
  \approx 4.47\times10^{-3}$ for \emph{every} $\alpha$ below the clamp;
  Eq.~\eqref{eq:arcsinh} tracks the true value to floating-point
  precision. Source:
  \code{notes/strengthening/C3\_arcsinh\_worked\_example\_results.md}.}
  \label{tab:arcsinh}
  \begin{tabular}{lccc}
    \toprule
    $\alpha$ & textbook arccosh & $\varepsilon$-clamp & arcsinh (Eq.~\ref{eq:arcsinh}) \\
    \midrule
    $10^{-4}$  & $1.00\times10^{-4}$            & $4.47\times10^{-3}$ & $1.00\times10^{-4}$ \\
    $10^{-6}$  & $1.00\times10^{-6}$            & $4.47\times10^{-3}$ & $1.00\times10^{-6}$ \\
    $10^{-8}$  & $0$ {\footnotesize(all bits lost)} & $4.47\times10^{-3}$ & $1.00\times10^{-8}$ \\
    $10^{-10}$ & $0$ {\footnotesize(all bits lost)} & $4.47\times10^{-3}$ & $1.00\times10^{-10}$ \\
    $0$ ($x=y$) & $0$ fwd, \code{NaN} bwd        & $4.47\times10^{-3}$ & $0$ fwd, $0$ bwd \\
    \bottomrule
  \end{tabular}
\end{table}

The $\varepsilon$-clamp removes the \code{NaN} at the price of a constant
forward bias of $\sim\!4.5\times10^{-3}$ and an extra hyperparameter;
Eq.~\eqref{eq:arcsinh} removes the \code{NaN} and keeps the value exact,
with one \code{safe\_sqrt} (a \code{torch.where} guarding the $\sqrt{0}$
derivative) as its only cost --- about $30\%$ per-call overhead in
isolation (batch $10^4$, $n = 8$), shrinking to a small fraction of any
realistic training step. The same reparameterization drives \code{log}.
The \code{safe\_*} idiom is folklore in \code{geoopt} and PyTorch3D
\citep{kochurov2020geoopt,ravi2020pytorch3d}; applied consistently across
the manifold API it keeps every boundary autograd-safe, and it reframes
the choice between the two distance forms as a precision decision.

\subsection{A reference-free correctness gate for operator-semigroup
primitives}
\label{sec:c7}

The bug in Section~\ref{sec:c1} was found by two checks that any correct
operator semigroup $S_t$ satisfies internally, each computable from the
object alone --- the only option for the higher-dimensional kernels,
where no reference solution existed:
\begin{description}\itemsep2pt
  \item[Operator-equation residual.] $S_t$ solves an evolution equation
  $\partial_t S = \mathcal{G} S$. Evaluate $\partial_t S - \mathcal{G} S$
  by finite differences and report it relative to $\lVert \mathcal{G}S\rVert$.
  A correct implementation sits at the finite-difference truncation
  floor.
  \item[Conserved-quantity check.] $S_t$ conserves a known quantity $Q$
  (probability mass, row sum, trace, energy). Evaluate $Q(S_t) - Q$;
  a correct implementation sits at machine precision.
\end{description}
Both are reference-free, which is what makes them usable for a
\emph{new} primitive. Applying them to a structurally different object
--- the graph heat kernel $K_t = e^{-tL}$ on a $24$-node graph, where the
PDE becomes the ODE $\dot K = -LK$ and probability mass becomes row mass
$K_t \mathbf{1} = \mathbf{1}$ (\code{notes/validation/graph\_heat\_kernel\_validation.py})
--- shows the two checks cover distinct failure modes
(Table~\ref{tab:dualcheck}).

\begin{table}[t]
  \centering
  \caption{Three corruptions of the graph heat kernel against the two
  checks. Each of the last two corruptions is invisible to one check and
  caught by the other; the two together form the gate.}
  \label{tab:dualcheck}
  \begin{tabular}{lcc}
    \toprule
    Corruption & Operator residual & Mass conservation \\
    \midrule
    spectral-shift $K \to e^{\alpha t} K$ & fires (grows with $t$) & fires ($=e^{\alpha t}-1$) \\
    constant rescale $K \to c\,K$         & \textbf{blind} (floor $\forall c$) & fires ($=|c-1|$) \\
    mass-preserving $K \to K + \varepsilon L'$ & fires ($\propto \varepsilon$) & \textbf{blind} (machine $\forall \varepsilon$) \\
    \bottomrule
  \end{tabular}
\end{table}

A constant rescale is in the null space of the residual --- a scalar
multiple of a solution to the homogeneous equation still solves it ---
but mass conservation catches it immediately. A mass-preserving
perturbation by another Laplacian $L'$ (so $L'\mathbf{1}=0$) leaves every
row sum unchanged but no longer solves the heat ODE, so the residual
fires. The spectral-shift bug of Section~\ref{sec:c1} happens to trip
both, which is why one check sufficed to \emph{find} it; the table shows
that in general both are needed. A closed-form spot check is a weaker
third gate, constraining only the points evaluated, and serves as
corroboration.

\section{Precision and speed}
\label{sec:precision}

\subsection{A closed form for the $\Hyp^5$ heat kernel}
\label{sec:c2}

The operator-chain form for odd $n = 2m+1$
\citep[Thm.~8.21]{grigoryan2009heat} is
\(
  p^{2m+1}(t,r) = \frac{(-1)^m}{(2\pi)^m}(4\pi t)^{-1/2}
  \big(\tfrac{1}{\sinh r}\partial_r\big)^m e^{-m^2 t - r^2/4t}.
\)
Carrying out the two operator applications for $m = 2$ by hand
(\code{notes/verification/heat\_kernel\_n5\_sympy.py}) gives the closed
form
\begin{equation}
  k^5_t(r) = (4\pi t)^{-5/2}\, e^{-4t - r^2/4t}\,
  \frac{r^2 \sinh r + 2t\,(r\cosh r - \sinh r)}{\sinh^3 r},
  \label{eq:n5}
\end{equation}
with the $r \to 0$ limit
$k^5_t(0) = (4\pi t)^{-5/2} e^{-4t}\,(1 + 2t/3)$. Symbolic verification
confirms that~\eqref{eq:n5} equals one corrected application
of~\eqref{eq:recursion} to $k^3$, satisfies the $\Hyp^5$ heat equation,
and matches the limit, all with residual exactly zero.

Using~\eqref{eq:n5} directly, rather than two autograd-based recursion
steps from $n = 3$, is both faster and more precise: each
\code{torch.autograd.grad} call accumulates floating-point noise, while
the closed form is one straight evaluation. The PDE residual drops by
about three orders of magnitude ($10^{-8}$--$10^{-6}$ versus
$10^{-5}$--$10^{-3}$). Table~\ref{tab:heatperf} reports latency: the
$n = 5$ closed form costs only slightly more than $n = 3$, an order of
magnitude below the even-$n$ quadrature-plus-recursion paths.

\begin{table}[t]
  \centering
  \caption{Heat-kernel latency, CPU \code{float64}, batch $4096$
  (\code{notes/strengthening/C1\_C2\_heat\_kernel\_bench\_results.md}).
  Commodity-GPU scaling is in Table~\ref{tab:gpuscale}
  (Section~\ref{sec:gpuscale}).}
  \label{tab:heatperf}
  \begin{tabular}{clrr}
    \toprule
    $n$ & path & forward (ms) & fwd+bwd (ms) \\
    \midrule
    3 & Davies--Mandouvalos closed form & 0.11 & 0.32 \\
    \textbf{5} & \textbf{hand-derived closed form (Eq.~\ref{eq:n5})} & \textbf{0.17} & \textbf{0.48} \\
    \textbf{7} & \textbf{hand-derived closed form (Eq.~\ref{eq:n7})} & \textbf{0.19} & \textbf{0.53} \\
    9 & recursion from $n=7$ (1 step) & 0.59 & 1.35 \\
    6 & recursion from $n=2$ (2 steps) & 11.2 & 21.7 \\
    \bottomrule
  \end{tabular}
\end{table}

\subsection{The same construction at $\Hyp^7$, and an independent check}
\label{sec:c2b}

One further operator step gives the $\Hyp^7$ closed form,
\begin{equation}
  k^7_t(r) = (4\pi t)^{-7/2}\, e^{-9t - r^2/4t}\,
  \frac{r^3\sinh^2 r + 6r^2 t\,\sinh r\cosh r
        + (8t^2 - 6t)\,r\sinh^2 r + 12t^2(r - \sinh r\cosh r)}{\sinh^5 r},
  \label{eq:n7}
\end{equation}
with $r \to 0$ limit $(4\pi t)^{-7/2} e^{-9t}(1 + 2t + 16t^2/15)$. It
equals one corrected application of~\eqref{eq:recursion} to $k^5$ and the
operator chain at $m = 3$ --- verified symbolically and numerically to
$\sim\!10^{-13}$, with an $\Hyp^7$ heat-equation residual $\sim\!10^{-12}$
(\code{notes/verification/heat\_kernel\_n7\_sympy.py}). Used directly it
costs essentially the same as the $n = 5$ closed form
(Table~\ref{tab:heatperf}), and it seeds the odd-$n$ recursion at
$n \ge 9$: one recursion step adds $\sim\!0.8$\,ms forward+backward
(compare $n = 9$ at $1.35$\,ms to the $n = 7$ closed form at
$0.53$\,ms), so the step's cost and its float-noise compounding are both
avoided at $n = 7$.

A Crank--Nicolson solver for the radial heat equation
$\partial_t u = \partial_r^2 u + (n-1)\coth(r)\,\partial_r u$, sharing
none of the operator-chain machinery, gives an independent confirmation:
evolving each library kernel forward in time reproduces it to the
scheme's truncation floor --- $7.8\times10^{-5}$ at $n = 5$,
$1.8\times10^{-4}$ at $n = 7$, $3.9\times10^{-4}$ at $n = 9$
(\code{notes/validation/heat\_kernel\_crank\_nicolson\_results.md}). A
wrong coefficient, the kind of error the spectral shift of
Section~\ref{sec:c1} introduced, would show an $O(1)$ mismatch rather
than $10^{-4}$.

\subsection{Letting a learnable curvature cross zero without retraining}
\label{sec:c4}

When $\kp$ is a learnable parameter on the $\kp$-stereographic model, the
geometry (spherical / hyperbolic / Euclidean) can change during
training. The map
\begin{equation}
  f(\kp, \alpha) = \frac{\tan_\kp^{-1}(\sqrt{|\kp|}\,\alpha)}{\sqrt{|\kp|}}
  = \begin{cases}
    \arctan(\sqrt{\kp}\,\alpha)/\sqrt{\kp}, & \kp > 0,\\
    \alpha, & \kp = 0,\\
    \operatorname{arctanh}(\sqrt{|\kp|}\,\alpha)/\sqrt{|\kp|}, & \kp < 0,
  \end{cases}
  \label{eq:fkappa}
\end{equation}
is a single analytic function of $\kp \in \R$, with Taylor series
$f(\kp,\alpha) = \alpha\sum_{m\ge0}(-\kp\alpha^2)^m/(2m+1)
= \alpha\,(1 - \kp\alpha^2/3 + \kp^2\alpha^4/5 - \cdots)$ converging for
$|\kp|\alpha^2 < 1$, exactly the manifold's distance-formula domain.
The two closed forms express one analytic function, one for each side of
zero. Symbolic verification
(\code{notes/verification/kappa\_crossing\_sympy.py}) confirms that the
two branches' values, first and second $\kp$-derivatives all agree at
$\kp = 0$, and that an integral representation
$f(\kp,\alpha)/\alpha = \int_0^1 (1+\kp(t\alpha)^2)^{-1}\,dt$ bridges both
signs without a conditional.

Realizing this through \code{torch.autograd} needs three pieces, since
autograd requires two finite expressions to mask between rather than one
analytic function: (1) clamp $|\kp|$ at \code{finfo.tiny} before the
square root, so $\sqrt{|\kp|}$ never vanishes and the scaled argument is
tiny-but-nonzero; (2) \code{safe} versions of $\arctan(t)/t$ and
$\operatorname{arctanh}(t)/t$ that short-circuit $t \le 0$ to the
analytic limit $1$; (3) an outer \code{torch.where(}$\kp>0$\code{, ...)}
selecting per element in forward and masking the gradient in backward.
Inside the manifold's distance-formula domain, both branches are finite,
so the dispatch is autograd-safe; outside it (an input the manifold's
own domain restriction excludes) the masked-out \code{arctanh} branch
can leak a \code{NaN} through the backward.

A $100$-step stochastic-gradient trajectory whose target oscillates
between $+0.7$ and $-0.7$ --- four sign crossings --- keeps $\kp$, its
gradient, the distances, and the points all finite at every step, with
no special-casing of $\kp = 0$
(\code{notes/strengthening/C4\_kappa\_crossing\_stress\_results.md}).
The alternatives are worse: a static-branch manifold that caches the
sign at construction computes $\operatorname{arctanh}$ where it should
compute $\arctan$ after a flip, a $1.38\%$ relative error ($0.4086$ vs.\
the correct $0.4030$); and a truncated unified Taylor series, the natural
\code{where}-free alternative, still leaves $\sim\!10^{-3}$ error with
$32$ terms near the domain boundary ($\sqrt{|\kp|}\alpha = 0.95$) where
the dispatch is exact. The dispatch evaluates both branches, but each is
a single elementwise kernel and the M\"obius machinery dominates and is
branch-independent, so the measured overhead is well below the
worst-case $2\times$: CPU $\sim\!1.2$--$1.4\times$ (batch $16384$,
$n=8$, \code{float64}), and on a consumer AMD Radeon RX 9060 XT (gfx1200,
ROCm 6.4, batch $65536$, $n=8$), $1.27\times$ in \code{float64} and
$1.58\times$ in \code{float32}. Curvature-as-learnable-scalar is
established \citep{bachmann2020constant,skopek2020mixed}; those
implementations leave the behavior \emph{at} $\kp = 0$ implicit,
sidestepping it by initializing on one side and regularizing away from
the boundary. The dispatch here crosses it, with the analyticity
verified symbolically and the crossing exercised in training.

\subsection{Scaling on a commodity GPU}
\label{sec:gpuscale}

The same kernels run at practical speed on consumer hardware. On an
RTX 3060 (12\,GB, \code{float64}, CUDA), forward+backward latency is
near-flat in batch while the same-machine CPU scales linearly: from a
batch of $4096$ to $\sim\!10^6$ (a $256\times$ increase) the GPU time
barely moves while the CPU time grows roughly proportionally
(Table~\ref{tab:gpuscale}). A batch of a million hyperbolic heat-kernel
evaluations costs on the GPU about what $4096$ does --- launch-bound, not
compute-bound, well past $10^6$ points --- and the CPU/GPU ratio there
reaches $\sim\!80$--$110\times$. With the $\kp$-crossing dispatch measured
on an AMD RX 9060 XT under ROCm (Section~\ref{sec:c4}), the substrate's
geometry runs on commodity cards of either vendor, no datacenter
accelerator required.

\begin{table}[t]
  \centering
  \caption{Heat-kernel forward+backward latency (ms), \code{float64},
  same-machine CPU vs.\ CUDA on an RTX 3060. GPU latency is near-flat in
  batch, CPU linear; the CPU/GPU ratio at $10^6$ is
  $\sim\!80$--$110\times$. Source:
  \code{notes/strengthening/gpu\_scaling\_results.md}.}
  \label{tab:gpuscale}
  \begin{tabular}{clrrrr}
    \toprule
    $n$ & path & CPU $4{\times}10^3$ & CPU $10^6$ & GPU $4{\times}10^3$ & GPU $10^6$ \\
    \midrule
    5 & closed form & 10.0 & 646 & 1.8 & 5.8 \\
    7 & closed form & 19.5 & 627 & 6.8 & 8.0 \\
    9 & recursion (from $n{=}7$) & 47 & 2320 & 9.7 & 24 \\
    \bottomrule
  \end{tabular}
\end{table}

\section{New primitives}
\label{sec:primitives}

\subsection{Continuous per-point curvature}
\label{sec:c5}

\code{HeterogeneousKappaManifold} exposes each point's curvature as a
continuous real-valued tensor: autograd flows through $\kp$ directly,
any $\kp \in \R$ is representable, and there is no expert bank to choose.
The closest prior art is discrete gating, in which each node mixes over a
fixed bank $\{\kp_1,\dots,\kp_K\}$ of constant-curvature spaces
\citep{guo2025graphmore}; its recovered curvature is confined to the
convex hull of the bank, so the bank must be guessed before training. The
class supplies the math --- per-point $\kp$-trig, exponential and
logarithm maps, and pair distance --- and leaves the $\kp$
parameterization to the caller. The $\kp$-field-plus-per-point-residual
parameterization used downstream belongs to the substrate project; the
library provides the geometry beneath it.

A controlled identifiability task makes the comparison concrete
(\code{notes/strengthening/C5\_C6\_synthetic\_task.py}): $90$ nodes in
three regions with true curvatures $\{-1.0,\,0.0,\,+0.5\}$, ground-truth
coordinates sampled per region, and pairwise distances under the true
$\kp$-field as embedding targets. Variants recover $(\text{coords},
\kp\text{-field})$ from distances (Table~\ref{tab:hetero}; mean $\pm$ std
over three seeds, $1500$ epochs, dimension $3$, on an AMD ROCm GPU).

\begin{table}[t]
  \centering
  \caption{Identifiability task: final fit loss and Pearson correlation
  between recovered and true per-node $\kp$. Continuous per-point $\kp$
  with the arithmetic-mean combiner attains the lowest loss and the best
  $\kp$-recovery; the adversarial-bank discrete model pays a
  bank-choice penalty of $\sim\!25\%$ in loss.}
  \label{tab:hetero}
  \begin{tabular}{llcc}
    \toprule
    Variant & combiner / bank & final loss & $\kp$-recovery corr. \\
    \midrule
    single $\kp$ baseline & --- & $0.0098 \pm 0.0046$ & N/A \\
    discrete gating & friendly $\{-1,0,+0.5\}$ & $0.0099 \pm 0.0042$ & $+0.630 \pm 0.071$ \\
    discrete gating & adversarial $\{-2,+1.5\}$ & $0.0124 \pm 0.0045$ & $+0.526 \pm 0.071$ \\
    \textbf{continuous per-point $\kp$} & \textbf{arithmetic\_mean} & $\mathbf{0.0065 \pm 0.0061}$ & $\mathbf{+0.667 \pm 0.077}$ \\
    continuous per-point $\kp$ & harmonic\_mean & $0.0092 \pm 0.0065$ & $-0.001 \pm 0.151$ \\
    continuous per-point $\kp$ & max\_magnitude & $0.0069 \pm 0.0055$ & $+0.250 \pm 0.198$ \\
    \bottomrule
  \end{tabular}
\end{table}

Continuous-with-arithmetic-mean wins on both metrics in the
truth-aligned case (the targets were generated under that combiner). The
single-$\kp$ baseline ties the friendly-bank model on loss; with one
degree of freedom it settles on a global $\kp \approx -0.84$ and
represents no per-node variation. Two caveats bound the result. The
$\kp$-recovery correlation tops out near $0.67$ because at dimension $3$
with $90$ nodes several $(\text{coords},\kp)$ configurations fit the
distances about equally well; the ranking across variants holds, while
the absolute number reflects the task's identifiability limit rather
than the primitive's capacity. And training continuous $\kp$ takes care:
a soft coordinate barrier ($\lVert x\rVert < 0.85$), a hard $\kp$ clamp
($|\kp| < 1.5$), gradient clipping, and cosine learning-rate decay keep
momentum from driving points outside the domain into a forward-\code{NaN}.
Autograd-safety governs the forward and backward passes; optimization
stability is a separate requirement, met here by those four guards.

\subsection{The pair-$\kp$ combiner abstraction}
\label{sec:c6}

Computing a pairwise distance between two points of different curvature
requires combining $\kp_x$ and $\kp_y$ into one effective pair-curvature,
and that rule remains a free choice in the literature. The library
exposes it as a callable $(\kp_x, \kp_y) \mapsto \kp_{\text{eff}}$ with
two built-in defaults and arbitrary override. Four combiners on the task above
(Table~\ref{tab:combiner}) show the design axis: the empirical winners
are smooth and well-defined at $\kp = 0$. \code{arithmetic\_mean} is both
the well-motivated default and the best performer; \code{harmonic\_mean}
collapses when the truth includes $\kp = 0$ because of the
\code{finfo.tiny} guard; the non-smooth combiners
(\code{signed\_geometric\_mean}'s $\operatorname{sign}$,
\code{max\_magnitude}'s $|\cdot|$) are brittle under Adam, and a
\code{safe\_sqrt} fixes the boundary derivative but not an upstream
non-smooth $\operatorname{sign}$. The abstraction's value is in surfacing
the choice: the combiner is a callable the user sets, with
\code{arithmetic\_mean} as a smooth, well-motivated default.

\begin{table}[t]
  \centering
  \caption{Combiner characterization. A combiner that is smooth and
  well-defined at $\kp=0$ trains reliably; a non-smooth one trains poorly.}
  \label{tab:combiner}
  \begin{tabular}{lccc}
    \toprule
    combiner & behavior at $\kp = 0$ & smoothness & fit on truth \\
    \midrule
    \code{arithmetic\_mean} (default) & Euclidean, well-defined & $C^\infty$ & best (matches truth) \\
    \code{harmonic\_mean} & needs \code{finfo.tiny} clamp & $C^\infty$ away from $0$ & collapses (clamp) \\
    \code{signed\_geometric\_mean} & $\operatorname{sign}$ jump at sum $=0$ & non-smooth & brittle across seeds \\
    \code{max\_magnitude} & jump at $|\kp_x|=|\kp_y|$ & non-smooth & mediocre \\
    \bottomrule
  \end{tabular}
\end{table}

\subsection{Mixed-curvature product manifolds}
\label{sec:c53}

For completeness the library includes \code{ProductManifold}, the
well-established construction of a product of constant-curvature factors
with a per-factor metric \citep{gu2019learning,skopek2020mixed}. The
projection and retraction compose factor-wise over the existing
single-manifold API. It reuses established geometry so the substrate
covers the mixed-curvature setting.

\subsection{Spectral dimension as a collapse diagnostic}
\label{sec:specdim}

A representation that nominally occupies $\R^D$ can collapse onto a far
lower-dimensional set; detecting that is one of the questions the
substrate program raises directly. The spectral dimension reads it off a
Laplacian spectrum: with return probability
$p(t) = \tfrac1n\sum_i e^{-t\lambda_i} \sim t^{-d_s/2}$, the estimator
$d_s = -2\,\mathrm{d}\log p / \mathrm{d}\log t$
(\code{spectral.spectral\_dimension}, after \citet{rammal1983}) recovers
a possibly non-integer dimension. Under one window rule --- the
asymptotic tail of the return probability --- it recovers integer lattice
dimensions and the Sierpinski gasket's fractional
$2\ln 3/\ln 5 \approx 1.365$ (Table~\ref{tab:specdim};
\code{notes/validation/spectral\_dimension\_results.md}). A point set
collapsed to an effective line registers $d_s \approx 1$ whatever $D$ is
--- the reading the substrate's observed drift toward low effective
dimension calls for.

\begin{table}[t]
  \centering
  \caption{\code{spectral\_dimension} against structures with known
  $d_s$, one window rule for all. The Sierpinski gasket is the canonical
  non-integer case \citep{rammal1983}; the lattice residuals are the
  finite-$t$ approach of the return-probability slope to its asymptote.}
  \label{tab:specdim}
  \begin{tabular}{lcc}
    \toprule
    structure & known $d_s$ & recovered \\
    \midrule
    1-D ring (cycle) & 1.000 & 1.000 \\
    2-D torus & 2.000 & 2.031 \\
    3-D torus & 3.000 & 3.094 \\
    Sierpinski gasket & 1.365 & 1.320 \\
    \bottomrule
  \end{tabular}
\end{table}

\section{Provenance: a content-addressable record for mechanistic interpretability}
\label{sec:provenance}

Studying whether and how geometric structure forms requires tracing which
primitive produced which representation, and intervening on it. The
library makes that native rather than an external instrumentation pass.
Every primitive decorated with \code{@with\_provenance} emits, inside a
\code{record()} context, a node in a content-addressable Merkle DAG keyed
by
\[
  \mathrm{hex}(\text{call}) = \mathrm{sha256}\big(\text{op\_id}
  + \text{version} + \mathrm{canonical}(\text{params})
  + [\,\mathrm{hex}(\text{input})\,]\big),
\]
leaf tensors hashed by content (bytes, shape, dtype). The same operation
on the same inputs yields the same hex, so identical sub-computations
across runs share a node. The DAG exports to NetworkX for graph
analysis, to a Pandas table of $(\text{op\_id}, \text{params},
\text{hex}, \text{shape})$ records for \mbox{SAELens}-style feature
training, and to JSON; substitution at a hex (\code{registry.substitute})
applies activation patching --- the TransformerLens intervention --- to
math primitives rather than network layers.

The scope is deliberate. This is a record at the \emph{math-primitive}
layer, complementary to neural-network hook libraries (TransformerLens,
nnsight) rather than a replacement; substitution does not auto-replay
downstream operations (the caller re-runs), and it is not an experiment
tracker. What it enables, by design, is a set of interpretability
operations expressed directly on the geometry: \emph{causal tracing} of a
representation back through the primitives that built it; \emph{ablation}
by substituting a node's output --- zeroing a curvature term, say --- and
re-running to measure the downstream effect; content-addressed
\emph{deduplication and caching} of repeated sub-computations; and
feature analysis over the exported record table. These are the
designed-for uses; demonstrating them at scale on a substrate model is
downstream work, not a result claimed here.

\section{Discussion}
\label{sec:discussion}

Across the seven results the mathematics was settled and the work lay in
the implementation --- making a finite-precision, automatically
differentiated, GPU-resident computation behave the way the mathematics
prescribes: a multiplicative factor autograd cannot see
(Section~\ref{sec:c1}), two algebraically identical distance formulas
with opposite numerical fates (Section~\ref{sec:c3}), a kernel that is
faster and more precise written out by hand than assembled by recursion
(Sections~\ref{sec:c2},~\ref{sec:c2b}), an analytic function that autograd can only reach
through a three-part safety net (Section~\ref{sec:c4}), and a per-point
curvature primitive whose autograd-safety is necessary but not sufficient
for stable optimization (Section~\ref{sec:c5}). The reference-free dual
check (Section~\ref{sec:c7}) certifies such a substrate where no
reference solution exists.

The bitter lesson predicts that hand-engineering gives way to scale and
general methods. This work acts on that prediction one level down. For
geometric structure to be learned directly instead of laundered through
language, the geometry it is learned on has to be correct, fast, and
differentiable, so that scale and search operate on the structure
itself. These results clear hand-engineering from the math layer --- the
per-project re-derivation, the silent precision loss, the boundary
\code{NaN}s, the curvature frozen at construction time --- and leave the
layer above to optimization. The delivered result is smaller and
checkable: a non-Euclidean embedding with learnable, per-point,
sign-crossing curvature stays exact and scales on commodity GPUs --- the
same video hardware the bitter lesson runs on, now carrying lossless
geometric structure rather than pixels. Curved embedding is, in this
concrete sense, practically viable and not only theoretically sound.
Whether structure made directly available shortens the route language
takes is the open wager Synoros was founded on; this work supplies the
substrate and leaves the verdict to the experiments it enables.

Three limits bound these claims. The empirical evidence is a controlled
identifiability task; large-scale downstream benchmarks belong to the
consuming project. The GPU figures come from two consumer cards, not a
datacenter accelerator or a full-scale model. And the emergence thesis of Section~\ref{sec:intro} motivates the work
while the seven substrate results are what it demonstrates --- they hold
whichever way the larger wager resolves. These are peer contributions to
an active direction
\citep{digiovanni2022heterogeneous,guo2025graphmore,gu2019learning}.
Natural next steps extend the substrate along the same line: a
closed-form $\Hyp^9$ kernel from the same operator chain, the
spectral-collapse diagnostic applied to a learned embedding rather than
known lattices, and downstream tasks that exercise per-point curvature at
scale.

\paragraph{Reproducibility.} Each result is backed by a symbolic
verification script (\code{notes/verification/}), a numerical
demonstration with recorded results (\code{notes/strengthening/},
\code{notes/validation/}), and a unit/property/comparison test suite in
\code{tests/}. Every numerical constant in the source is derived, marked
as a universal invariant, or cataloged in \code{notes/magic\_numbers.md},
and a static audit enforces this in continuous integration.

\bibliographystyle{plainnat}
\bibliography{references}

\end{document}
